% Generated draft: requires mathematical and source review.
\documentclass[11pt]{article}
\usepackage[margin=25mm]{geometry}
\usepackage{fontspec}
\setmainfont{Latin Modern Roman}
\usepackage{amsmath,amssymb,mathtools,unicode-math}
\setmathfont{Latin Modern Math}
\usepackage{enumitem,needspace,xurl,hyperref}
\setlength{\parindent}{0pt}
\setlength{\parskip}{5pt}
\setlength{\emergencystretch}{3em}
\newcommand{\N}{\mathbb{N}}
\newcommand{\Z}{\mathbb{Z}}
\newcommand{\Q}{\mathbb{Q}}
\newcommand{\R}{\mathbb{R}}
\newcommand{\C}{\mathbb{C}}
\newcommand{\F}{\mathbb{F}}
\newcommand{\A}{\mathbb{A}}
\newcommand{\PP}{\mathbb{P}}
\newcommand{\E}{\mathbb{E}}
\newcommand{\eps}{\varepsilon}
\newcommand{\dd}{\,\mathrm{d}}
\begin{document}
\section*{UnsolvedMath ID: 10600043}
\textbf{Source problem code:} catalogue entry
\subsection{Definitions and mathematical statement}
The source describes Sarkar's construction, initially titled ``Knot Floer Homotopy,'' of a cell complex whose homology agrees with Floer homology, using the Manolescu--Ozsv\'ath--Sarkar--Szab\'o--Thurston grid-diagram approach and shellability. It asks for an analogous construction for Khovanov homology, suggesting Bloom's construction as a more promising starting point than ordering the chains/states in the original Khovanov definition.

For a diagram \(D\), write \(\operatorname{Kh}^{i,j}(D;A)\) for the intended Khovanov groups, where \(i\) is homological and \(j\) quantum grading and \(A\) is the chosen coefficient ring. A precise homological realization might seek combinatorially defined based cell or simplicial complexes \(X_j(D)\), with stated shifts \(s_j\), satisfying
\[\widetilde H_{i+s_j}(X_j(D);A)\cong\operatorname{Kh}^{i,j}(D;A),\]
or a specified equivalent realization appropriate to the chosen grading convention. The shifts/reduced conventions are necessary formulation choices rather than a claim that ordinary unshifted homology realizes negative chain degrees. If a cohomological or stable-spectrum realization is intended instead, its relationship to the source's homology/cell-complex wording must be explained rather than silently substituted.

The construction should specify the combinatorial cells or simplices and attachment/incidence data, identify the resulting chain theory with the chosen Khovanov theory, and state the resulting diagram-invariance property at the homology and, if claimed, homotopy/stable levels. The excerpt gives no actual ordering, poset, shelling or attachment rule and does not assert that every chain complex automatically has such a realization. The precise Bloom construction, Khovanov version and scope of diagrams are absent. Although the record belongs to a virtual-knot survey and the research summary mentions all virtuals, the quoted request itself does not explicitly restrict to or quantify over all virtual links; that extension requires a specified virtual Khovanov convention.
\subsection{Short English statement}
Construct a combinatorial cell or simplicial realization of Khovanov homology analogous to Sarkar's Floer construction, potentially using Bloom's model. The grading, homology-versus-cohomology and classical-versus-virtual scope must be fixed.
\subsection{Sources}
\begin{itemize}
\item Local UnsolvedMath registry, permanent ID \(10600043\), exact source code \texttt{catalogue entry}.
\item \emph{Unsolved Problems in Virtual Knot Theory and Combinatorial Knot Theory}, revised survey, problem-list item 43: \url{https://arxiv.org/abs/1409.2823}. Registry access 2026-07-29; source-tex extraction. Embedded citations for Sarkar and Bloom are elided. ``Knot Floer Homotopy'' is reported as an initial title, not a verified final bibliographic title.
\item The source names the Manolescu--Ozsv\'ath--Sarkar--Szab\'o--Thurston grid approach. The research summary attributes substantial analogous progress to Lipshitz--Sarkar and mentions thickened-surface extensions, without precise citations.
\end{itemize}
\end{document}
